\documentclass[11pt]{article}
\usepackage[a4paper,margin=25mm]{geometry}
\usepackage{amsmath,amssymb,amsthm,mathtools,hyperref}
\hypersetup{hidelinks}
\emergencystretch=2em
\newcommand{\E}{\mathbb E}
\newcommand{\Pp}{\mathbb P}
\newcommand{\R}{\mathbb R}
\newcommand{\norm}[1]{\lVert#1\rVert}
\newtheorem{theorem}{Theorem}
\newtheorem{lemma}[theorem]{Lemma}
\newtheorem{proposition}[theorem]{Proposition}
\newtheorem{corollary}[theorem]{Corollary}
\theoremstyle{remark}
\newtheorem*{remark}{Remark}
\title{Bounds for all recursive IRMF components\\
under a curvature-dispersion condition}
\author{}
\date{31 October 2026}
\begin{document}
\maketitle

\noindent This note proves the propagation remainder, closes the induction,
and compares the fixed-noise reference with the paper's population targets.
The proof adds an explicit curvature-dispersion condition to the assumptions
in the submitted draft. It also states a quantitative finest-scale condition.
The population-target comparison additionally uses a fourth noise moment.
These are sufficient conditions, not claims that the original assumptions
alone imply the result. No bound on the unknown propagation remainder is
assumed.

\section{Setting and assumptions}
Let \(Y=X+\varepsilon\in\R^n\), where \(X\) is deterministic and the
\(\varepsilon_u\) are independent with a common law. Use the
Gaussian-smoothed absolute loss \(\rho=\rho_H\), \(H>0\), and constants
\[
 |\rho'|\le1,\qquad 0<\rho''\le L_2,\qquad |\rho'''|\le L_3,
 \quad L_2=\frac{\sqrt{2/\pi}}H,\quad
 L_3=\frac{\sqrt{2/\pi}e^{-1/2}}{H^2}.
\]
To allow laws without a first moment, define
\(\phi(c)=\E[\rho(c+\varepsilon)-\rho(\varepsilon)]\).
This is finite because the integrand is bounded by \(|c|\).
When the expected loss is finite, its minimizers are unchanged.

Let \(a_{tu}^{(k)}\ge0\), \(\sum_u a_{tu}^{(k)}=1\), be deterministic
weights, \(1\le k\le K\), and let \(\alpha_k=\max_{t,u}a_{tu}^{(k)}\).
Write
\[
 \mathcal M_k(y)_t=\arg\min_s\sum_u a_{tu}^{(k)}\rho(y_u-s),
 \quad Y^{(1)}=Y,\quad
 \widetilde S^{(k)}=\mathcal M_k(Y^{(k)}),\quad
 Y^{(k+1)}=Y^{(k)}-\widetilde S^{(k)}.
\]
The deterministic reference is
\[
 X_\phi^{(1)}=X,\quad
 S_{\phi,t}^{(k)}=\arg\min_s\sum_u a_{tu}^{(k)}\phi(X_{\phi,u}^{(k)}-s),
 \quad X_\phi^{(k+1)}=X_\phi^{(k)}-S_\phi^{(k)}.
\]
These minimizers exist and are unique: the scores are strictly increasing
in their argument and tend to \(\pm1\).
For \(c_{tu}^{(k)}=X_{\phi,u}^{(k)}-S_{\phi,t}^{(k)}\), define
\begin{align*}
 g_{ktu}(\varepsilon_u)&=\rho''(c_{tu}^{(k)}+\varepsilon_u),&
 f_{kt}&=\sum_u a_{tu}^{(k)}\E g_{ktu},\\
 q_{kt}&=\sum_u a_{tu}^{(k)}\rho'(c_{tu}^{(k)}+\varepsilon_u),&
 b_{kt}&=q_{kt}/f_{kt},&
 V_k[t,u]&=a_{tu}^{(k)}\E g_{ktu}/f_{kt}.
\end{align*}
The population score equation gives \(\E q_{kt}=0\).
Every \(V_k\) is deterministic and row stochastic. Throughout, assume
\begin{align}
 f_{kt}&\ge\lambda>0,\label{eq:curvature}\\
 \alpha_j&\le C_\alpha\alpha_k a^{-(k-j)}
       \quad(j\le k),\qquad a>1,\quad C_\alpha\ge1,\label{eq:schedule}\\
 \norm{B_{m:l}}_\infty&\le C_B\quad
       (1\le l\le m\le K),\qquad C_B\ge1,\label{eq:products}\\
 B_{m:l}&=(I-V_m)\cdots(I-V_l),\qquad B_{m:l}=I\quad(m<l).\notag
\end{align}
Here the matrix norm is the maximum absolute row sum. The verified Fourier
and perturbation lemmas in \texttt{irmf\_recursive\_component\_bounds.tex}
give one route to \eqref{eq:products}; a separate deterministic bound on
these actual population products is also sufficient.

The additional assumption concerns the known, single-observation noise law:
\begin{equation}
 d:=\max_{k,t}\frac{\sum_u a_{tu}^{(k)}
                \E|g_{ktu}-\E g_{ktu}|}{f_{kt}}
       <\frac{a-1}{C_BC_\alpha}.
 \label{eq:dispersion}
\end{equation}
This is a bound on relative curvature variation, not a bound on an unknown
recursive error. It permits contamination if the resulting law satisfies it.
The proof below does not cover all contamination laws at a fixed \(H\).

\section{Exact decomposition}
Set
\[
 e_k=\widetilde S^{(k)}-S_\phi^{(k)},\quad
 E_k=\sum_{j<k}e_j,\quad
 R_k=\mathcal M_k(X_\phi^{(k)}+\varepsilon)-S_\phi^{(k)}-b_k,
\]
and
\[
 \Phi_k=\mathcal M_k(X_\phi^{(k)}+\varepsilon)
              -\mathcal M_k(Y^{(k)})-V_kE_k.
\]
Then \(Y^{(k)}=X_\phi^{(k)}+\varepsilon-E_k\) and
\begin{equation}
 E_{k+1}=(I-V_k)E_k+b_k+R_k-\Phi_k. \label{eq:recursion}
\end{equation}
Define deterministic-transfer leading terms and accumulated remainders by
\begin{align}
 \Lambda_1&=\Theta_1=0,&
 \Lambda_{k+1}&=(I-V_k)\Lambda_k+b_k,&
 \Theta_k&=E_k-\Lambda_k, \label{eq:lambda}\\
 L_k&=b_k-V_k\Lambda_k,&
 Q_k&=e_k-L_k=\Theta_{k+1}-\Theta_k. \notag
\end{align}
In particular,
\begin{equation}
 \Lambda_k=\sum_{j<k}B_{k-1:j+1}b_j,\qquad
 \Theta_k=\sum_{j<k}B_{k-1:j+1}(R_j-\Phi_j).
 \label{eq:unroll}
\end{equation}
All stages use the same original noise vector. Independence is used only
across its coordinates.

\section{A quadratic-statistic bound with a complete proof}
\begin{lemma}\label{lem:chaos}
Let \(Z_1,\ldots,Z_n\) be independent. For \(u\ne v\), suppose
\(h_{uv}(Z_u,Z_v)\) is centered in each argument separately and
\(|h_{uv}|\le H_{uv}\), with deterministic \(H_{uv}\).
For every even integer \(p\ge2\),
\[
 \left\|\sum_{u\ne v}h_{uv}(Z_u,Z_v)\right\|_{L^p}
       \le16p\left(\sum_{u\ne v}H_{uv}^2\right)^{1/2}.
\]
Consequently, for \(z\ge2\), outside a set of probability \(e^{-z}\),
\[
 \left|\sum_{u\ne v}h_{uv}(Z_u,Z_v)\right|
       \le32\mathrm e\,z\left(\sum_{u\ne v}H_{uv}^2\right)^{1/2}.
\]
The kernels need not be symmetric or identical.
\end{lemma}
\begin{proof}
Choose a random subset \(A\) by independent fair coins, independently of
the \(Z_u\). Each ordered pair \(u\ne v\) belongs to \(A\times A^c\) with
probability \(1/4\). Thus the full sum is four times the conditional
average of its cross-partition sum. Jensen's inequality bounds its
\(L^p\) norm by four times the largest \(L^p\) norm of such a sum.

Fix a partition. Conditional on the variables in \(A\), symmetrize the
independent centered contributions indexed by \(v\in A^c\). This costs
a factor two and introduces independent Rademacher signs \(s_v\).
Conditional on \(Z_{A^c}\) and these signs, symmetrize over \(u\in A\).
Centering in the other argument justifies this second step, which costs
another factor two and introduces signs \(r_u\).
The usual symmetrization inequality follows directly by subtracting an
independent copy, applying conditional Jensen, inserting signs into the
symmetric differences, and applying the triangle inequality.

For even \(p\), the Rademacher moment inequality
\[
 \left\|\sum_i r_i a_i\right\|_{L^p}
       \le\sqrt p\left(\sum_i a_i^2\right)^{1/2}
\]
follows by expanding the moment and comparing even multiplicities with a
Gaussian moment, bounded by \(p^{p/2}\).
Apply it twice, with Minkowski's inequality in \(L^{p/2}\) between
applications, to obtain
\[
 \left\|\sum_{u\in A,v\in A^c}r_us_v h_{uv}(Z_u,Z_v)
        \right\|_{L^p(r,s)}
 \le p\left(\sum_{u\in A,v\in A^c}|h_{uv}(Z_u,Z_v)|^2\right)^{1/2}.
\]
The envelope makes the right side deterministic. The factors are
\(4\cdot2\cdot2=16\). Finally take \(p=2\lceil z/2\rceil\), so
\(z\le p\le2z\), and apply Markov's inequality at
\(\mathrm e\) times the \(L^p\) bound.
\end{proof}

\section{Constants and one simultaneous event}
Define
\begin{align*}
 c_u&=\sqrt2+2/3,&c_v&=\sqrt2+1/3,&
 s_1&=\frac{C_\alpha}{1-a^{-1}},&
 s_2&=\frac{\sqrt{C_\alpha}}{1-a^{-1/2}},&
 r&=\frac{C_\alpha}{a},\\
 A_\Lambda&=\frac{C_B}{\lambda}(\sqrt2s_2+2s_1/3),&
 G_\infty&=\frac{2C_Bs_1}{\lambda},&
 G_2&=\frac{2C_Bs_2}{\lambda},&
 b_0&=\frac{L_2}{\lambda},\\
 A_R&=\frac{2L_2c_uc_v}{\lambda^2}
              +\frac{2L_3c_u^2}{\lambda^3},&
 C_{\rm loc}&=\frac{L_2c_v}{\lambda}+\frac{4L_3c_u}{\lambda^2}.&&
\end{align*}
For \(x>0\), let \(z_0=\log(10nK)\) and \(z=z_0+x\).
The operator of centered curvature fluctuations is
\[
 Z_k[t,u]=\frac{a_{tu}^{(k)}}{f_{kt}}(g_{ktu}-\E g_{ktu}).
\]
This notation \(Z_k\) is a random matrix, not a fresh noise vector.

\begin{lemma}\label{lem:event}
If \(z\bar\alpha\le1\), where \(\bar\alpha=\max_k\alpha_k\), there is
one event of probability at least \(1-e^{-x}\) on which all the following
bounds hold simultaneously:
\begin{align}
 |q_{kt}|&\le c_u\sqrt{z\alpha_k},&
 \left|\sum_u a_{tu}^{(k)}g_{ktu}-f_{kt}\right|
       &\le L_2c_v\sqrt{z\alpha_k},\label{eq:score-hess}\\
 \norm{Z_k}_\infty&\le d+b_0c_v\sqrt{z\alpha_k},\label{eq:Znorm}\\
 \norm{\Lambda_k}_\infty&\le A_\Lambda\sqrt{z\alpha_{k-1}},
       \quad k=2,\ldots,K+1,\label{eq:Lambda-bound}\\
 \norm{L_k}_\infty&\le A_\Lambda\sqrt{z\alpha_k},\label{eq:L-bound}\\
 \norm{Z_k\Lambda_k}_\infty
   &\le b_0G_\infty\alpha_{k-1}
          +32\mathrm e\,b_0G_2 z\sqrt{\alpha_k\alpha_{k-1}},
          \quad k\ge2.\label{eq:ZLambda}
\end{align}
\end{lemma}
\begin{proof}
The score and curvature estimates follow from two-sided Bernstein bounds
using \(\sum_u(a_{tu}^{(k)})^2\le\alpha_k\), centered summand bounds
\(2a_{tu}^{(k)}\) and \(L_2a_{tu}^{(k)}\), and variance proxies
\(\alpha_k\) and \(L_2^2\alpha_k\). Their union costs at most
\(4nKe^{-z}\). For \eqref{eq:Znorm}, apply the one-sided bound to the
independent variables \(|g_{ktu}-\E g_{ktu}|/f_{kt}\in[0,b_0]\).
Their mean weighted sum is at most \(d\), and the union costs \(nKe^{-z}\).

To bound \(\Lambda_k\), group its representation \eqref{eq:unroll} by
observation. An uncentered envelope for the stage-\(j\) contribution in
one row has \(\ell_1\) norm at most \(C_B/\lambda\), maximum at most
\(C_B\alpha_j/\lambda\), and \(\ell_2\) norm at most
\(C_B\sqrt{\alpha_j}/\lambda\).
Thus its sum has variance proxy
\((C_B/\lambda)^2(\sum_{j<k}\sqrt{\alpha_j})^2\) and centered increment
bound \(2(C_B/\lambda)\sum_{j<k}\alpha_j\).
Equation \eqref{eq:schedule} bounds these two sums by
\(s_2\sqrt{\alpha_{k-1}}\) and \(s_1\alpha_{k-1}\).
Bernstein and \(z\alpha_{k-1}\le1\) give \eqref{eq:Lambda-bound}.
The union over \(k=2,\ldots,K+1\) costs \(2nKe^{-z}\).
For \(L_k\), its transfer operators are \(I\) at \(j=k\) and
\(-V_kB_{k-1:j+1}\) at \(j<k\). Their norms are at most \(C_B\)
because \(\norm{V_k}_\infty=1\).
The same argument, with sums through \(k\), gives \eqref{eq:L-bound}
and costs another \(2nKe^{-z}\).

For the last estimate, center each observation's score contribution:
\[
 \Lambda_k(u)=\sum_v\gamma_{uv}^{(k)}(\varepsilon_v),\qquad
 \E\gamma_{uv}^{(k)}=0,\qquad
 |\gamma_{uv}^{(k)}|\le\Gamma_{uv}^{(k)},
\]
where a valid deterministic envelope is
\[
 \Gamma_{uv}^{(k)}
 =\frac2\lambda\sum_{j<k}\sum_s
       |B_{k-1:j+1}[u,s]|a_{sv}^{(j)}.
\]
Minkowski's inequality and \eqref{eq:schedule} give
\[
 \max_{u,v}\Gamma_{uv}^{(k)}\le G_\infty\alpha_{k-1},\qquad
 \max_u\left(\sum_v(\Gamma_{uv}^{(k)})^2\right)^{1/2}
       \le G_2\sqrt{\alpha_{k-1}}.
\]
For fixed \(k,t\), expand \( (Z_k\Lambda_k)_t
=\sum_{u,v}Z_k[t,u]\gamma_{uv}^{(k)}(\varepsilon_v)\).
The diagonal is bounded in absolute value by
\(b_0G_\infty\alpha_{k-1}\).
Off the diagonal the terms are centered in each independent argument.
Their envelopes \(b_0a_{tu}^{(k)}\Gamma_{uv}^{(k)}\) have squared sum at
most \(b_0^2G_2^2\alpha_k\alpha_{k-1}\).
Lemma~\ref{lem:chaos} and a union bound cost at most \(nKe^{-z}\).
All event costs sum to \(10nKe^{-z}=e^{-x}\).
\end{proof}

\section{Propagation estimate}
\begin{lemma}\label{lem:propagation}
Suppose \eqref{eq:score-hess} holds and
\[
 L_2c_v\sqrt{z\alpha_k}\le\lambda/4,\qquad
 c_u\sqrt{z\alpha_k}\le\lambda^2/(8L_3).
\]
Then \(\norm{R_k}_\infty\le A_Rz\alpha_k\). Furthermore, for every
realized \(E_k\), with \(\Theta_k=E_k-\Lambda_k\),
\begin{align}
 \norm{\Phi_k}_\infty
 \le {}&\frac{4L_3}{\lambda}\norm{E_k}_\infty^2
       +C_{\rm loc}\sqrt{z\alpha_k}\norm{E_k}_\infty \notag\\
      &+\norm{Z_k\Lambda_k}_\infty
       +\norm{Z_k}_\infty\norm{\Theta_k}_\infty .
 \label{eq:propagation}
\end{align}
No independence between \(E_k\) and the noise is assumed.
\end{lemma}
\begin{proof}
Fix a row and let \(y=X_\phi^{(k)}+\varepsilon\).
At the population target the empirical curvature differs from \(f_{kt}\)
by at most \(v=L_2c_v\sqrt{z\alpha_k}\), and the score is at most
\(u=c_u\sqrt{z\alpha_k}\).
Within distance \(2u/\lambda\), the curvature is at least
\(\lambda-v-2L_3u/\lambda\ge\lambda/2\).
The empirical score therefore has its zero in this interval.
If its displacement is \(h\), Taylor's formula yields
\[
 |h-q_{kt}/f_{kt}|
 \le\frac{v|h|+L_3h^2/2}{\lambda}
 \le\frac{2uv}{\lambda^2}+\frac{2L_3u^2}{\lambda^3}
 \le A_Rz\alpha_k .
\]

The local fit is smooth, coordinatewise monotone, and translation
equivariant. Hence it is 1-Lipschitz in the sup norm.
Its Jacobian at an input \(w\) has row
\[
 J(w)[t,u]=
 \frac{a_{tu}^{(k)}\rho''(w_u-\mathcal M_k(w)_t)}
      {\sum_v a_{tv}^{(k)}\rho''(w_v-\mathcal M_k(w)_t)}.
\]
It is nonnegative with row sum one. For nonnegative vectors \(p,q\) of
positive mass,
\begin{equation}
 \left\|\frac p{\sum p}-\frac q{\sum q}\right\|_1
       \le\frac{2\|p-q\|_1}{\sum q}.
 \label{eq:normalization}
\end{equation}
This follows by adding and subtracting \(p/\sum q\).

The empirical curvature at the single-pass fit is at least \(\lambda/2\).
Let \(J(s)=J(y-sE_k)\). The Lipschitz property gives
\[
 |(y_u-sE_k(u)-\mathcal M_k(y-sE_k)_t)
       -(y_u-\mathcal M_k(y)_t)|
       \le2s\norm{E_k}_\infty .
\]
Apply \eqref{eq:normalization}, using the curvature at \(s=0\) in its
denominator, to get
\(\norm{J(s)-J(0)}_\infty\le8L_3s\norm{E_k}_\infty/\lambda\).
Thus the integrated quadratic error is at most
\(4L_3\norm{E_k}_\infty^2/\lambda\).
Only positivity, not a uniform lower curvature bound, is needed at
the other points of the segment.

Set \(\overline M[t,u]=a_{tu}^{(k)}g_{ktu}/f_{kt}=V_k[t,u]+Z_k[t,u]\).
Adding and subtracting the fitted-curvature numerator divided by \(f_{kt}\)
gives
\[
 \norm{J(0)-\overline M}_\infty
 \le\frac{v+4L_3u/\lambda}{\lambda}
 \le C_{\rm loc}\sqrt{z\alpha_k}.
\]
Finally, the exact integral identity is
\[
 \Phi_k=\int_0^1(J(s)-J(0))E_k\,ds
          +(J(0)-\overline M)E_k+Z_k\Lambda_k+Z_k\Theta_k .
\]
Taking norms proves \eqref{eq:propagation}.
\end{proof}

\section{Closing the induction with explicit constants}
Choose any fixed
\[
 d<\delta<\frac{a-1}{C_BC_\alpha},\qquad
 q=\frac{C_BC_\alpha\delta}{a-1}<1,\qquad A=2A_\Lambda .
\]
The midpoint of the displayed interval is one possible choice of \(\delta\).
Put
\begin{align}
 P(A)&=\frac{4L_3}{\lambda}A^2r+C_{\rm loc}A\sqrt r
          +b_0G_\infty r+32\mathrm e\,b_0G_2\sqrt r,\label{eq:P}\\
 D&=\frac{C_Bs_1\{A_R+P(A)\}}{1-q},\label{eq:D}\\
 \eta&=\min\left\{1,\
   \left(\frac{\lambda}{4L_2c_v}\right)^2,\
   \left(\frac{\lambda^2}{8L_3c_u}\right)^2,\
   \left(\frac{\delta-d}{b_0c_v}\right)^2,\
   \left(\frac{A_\Lambda}{D}\right)^2\right\}.\label{eq:eta}
\end{align}
Every quantity here is deterministic and specified before the induction.

\begin{theorem}[All components relative to the fixed-noise path]\label{thm:main}
Under \eqref{eq:curvature}--\eqref{eq:dispersion}, let \(x>0\) satisfy
\((x+\log(10nK))\bar\alpha\le\eta\).
With probability at least \(1-e^{-x}\), simultaneously for every \(k\le K\),
\begin{align}
 \norm{\Theta_{k+1}}_\infty&\le Dz\alpha_k,&
 \norm{E_{k+1}}_\infty&\le A\sqrt{z\alpha_k},\label{eq:closed}\\
 \norm{\Phi_k}_\infty&\le\{P(A)+\delta Dr\}z\alpha_k,\label{eq:Phi-final}\\
 \norm{\widetilde S^{(k)}-S_\phi^{(k)}-L_k}_\infty
       &\le D(1+r)z\alpha_k,\label{eq:expansion}\\
 \norm{\widetilde S^{(k)}-S_\phi^{(k)}}_\infty
       &\le A_\Lambda\sqrt{z\alpha_k}+D(1+r)z\alpha_k.\label{eq:main}
\end{align}
Here \(\Phi_1=0\). The constants do not depend on \(k\); they remain uniform
over \(K\) if the assumed parameters and the dispersion margin do.
\end{theorem}
\begin{proof}
Work on the single event in Lemma~\ref{lem:event}.
Equation \eqref{eq:eta} ensures the single-pass conditions and
\(\norm{Z_j}_\infty\le\delta\).
The base case is \(E_1=\Lambda_1=\Theta_1=0\).

Suppose \(\norm{\Theta_j}_\infty\le Dz\alpha_{j-1}\) for all \(2\le j\le k\).
Then \eqref{eq:Lambda-bound} and \(D\sqrt{z\bar\alpha}\le A_\Lambda\)
give \(\norm{E_j}_\infty\le A\sqrt{z\alpha_{j-1}}\).
Since \(\alpha_{j-1}\le r\alpha_j\) and \(z\ge1\), insert these bounds
and \eqref{eq:ZLambda} into \eqref{eq:propagation} to obtain
\[
 \norm{\Phi_j}_\infty
       \le P(A)z\alpha_j+\delta Dz\alpha_{j-1}.
\]
For \(j=1\) use \(\Phi_1=0\) and set \(\alpha_0=0\).
Unrolling \(\Theta\), rather than multiplying one-step norm bounds, yields
\begin{align*}
 \norm{\Theta_{k+1}}_\infty
 &\le C_Bz\left\{(A_R+P(A))\sum_{j=1}^k\alpha_j
                   +\delta D\sum_{j=1}^{k-1}\alpha_j\right\}\\
 &\le\left\{C_Bs_1(A_R+P(A))
              +\frac{C_BC_\alpha\delta}{a-1}D\right\}z\alpha_k
 =Dz\alpha_k .
\end{align*}
This closes the induction because \eqref{eq:D} solves the displayed
inequality with \(q<1\), while \eqref{eq:eta} fixed \(A\) in advance.
It also proves the \(E\) bound at the next stage.
Finally \(Q_k=\Theta_{k+1}-\Theta_k\), so
\(\norm{Q_k}_\infty\le D(1+r)z\alpha_k\); combine with
\eqref{eq:L-bound}.
\end{proof}

\begin{remark}
If adjacent weights also satisfy \(\alpha_k\le a_+\alpha_{k-1}\), then
\eqref{eq:Phi-final} is a bound of order \(z\alpha_{k-1}\), as in the
original draft. The theorem does not require this extra upper ratio.
The quadratic-statistic constant and the use of sup norms are conservative;
the theorem is not a numerical certificate for every small-window experiment.
\end{remark}

\section{Comparison with the paper's population targets}
Assume now \(\mu_4=\E|\varepsilon|^4<\infty\), and define
\[
 S_{{\rm pdf},t}^{(k)}
       =\arg\min_s\sum_u a_{tu}^{(k)}\E\rho(Y_u^{(k)}-s).
\]
Let \(m_\rho\) be the unique root of
\(\E\rho'(\varepsilon-m_\rho)=0\).
The paper's target at each stage depends only on the law of the actual
residual \(Y^{(k)}\); its auxiliary clean recursion need not be substituted
for \(X_\phi^{(k)}\).

Suppose
\[
 x_*=\eta/\bar\alpha-z_0>0,\qquad
 M_4=2^{K+1}\{\norm X_\infty+(n\mu_4)^{1/4}+|m_\rho|\}.
\]
For \(k\ge2\) set
\begin{align}
 U_k&=D\alpha_{k-1}(z_0+1)
       +(M_4+C_BK/\lambda)e^{-x_*/2},\label{eq:U}\\
 V_k^{\rm mom}&=A^2\alpha_{k-1}(z_0+1)+M_4^2e^{-x_*/2},\label{eq:Vmom}\\
 \beta_k&=L_2G_\infty\alpha_{k-1}
                  +L_2U_k+(L_3/2)V_k^{\rm mom}.\label{eq:beta}
\end{align}
The matrix \(V_k\) and scalar \(V_k^{\rm mom}\) are different quantities.

\begin{proposition}[Deterministic target correction]\label{prop:targets}
If, for \(2\le k\le K\),
\begin{equation}
 L_3\sqrt{V_k^{\rm mom}}\le\lambda/4,\qquad
 \beta_k\le\lambda^2/(8L_3),\label{eq:target-smallness}
\end{equation}
then
\[
 \norm{S_{\rm pdf}^{(k)}-S_\phi^{(k)}}_\infty
       \le2\beta_k/\lambda.
\]
At \(k=1\) the two targets coincide.
\end{proposition}
\begin{proof}
First, the high-probability estimates imply
\begin{equation}
 \E\norm{\Theta_k}_\infty\le U_k,\qquad
 \E\norm{E_k}_\infty^2\le V_k^{\rm mom}.
 \label{eq:moments}
\end{equation}
Here are the needed tail details. An empirical fit lies between its
smallest and largest input. A \(\phi\)-fit lies between
\(\min_u x_u+m_\rho\) and \(\max_u x_u+m_\rho\).
Consequently, for all relevant stages,
\[
 \norm{E_k}_\infty
 \le2^{K+1}(\norm X_\infty+\norm\varepsilon_\infty+|m_\rho|),
 \qquad \norm{\Lambda_k}_\infty\le C_BK/\lambda.
\]
Since \(\E\norm\varepsilon_\infty^4\le n\mu_4\), the \(L^4\) norm of
the first bound is at most \(M_4\). Thus the \(L^2\) norm of
\(\norm{\Theta_k}_\infty\) is at most \(M_4+C_BK/\lambda\).

Integrate the tails from Theorem~\ref{thm:main} for \(0<x\le x_*\).
For \(\Theta_k\), integration up to \(D\alpha_{k-1}(z_0+x_*)\)
gives at most \(D\alpha_{k-1}(z_0+1)\). Beyond this threshold, the
event is contained in the exceptional event at \(x_*\).
Cauchy--Schwarz bounds its contribution by
\((M_4+C_BK/\lambda)e^{-x_*/2}\).
Apply the same argument to \(\norm{E_k}_\infty^2\), whose \(L^2\) norm
is at most \(M_4^2\), to obtain \eqref{eq:moments}.

At \(S_{\phi,t}^{(k)}\), the paper's population score equals
\(\sum_u a_{tu}^{(k)}\E\rho'(c_{tu}^{(k)}+\varepsilon_u-E_k(u))\).
Taylor expansion around \(c_{tu}^{(k)}+\varepsilon_u\) gives
\[
 \left|\text{population score}\right|
 \le\left|\sum_u a_{tu}^{(k)}\E[g_{ktu}\Lambda_k(u)]\right|
       +L_2\E\norm{\Theta_k}_\infty
       +(L_3/2)\E\norm{E_k}_\infty^2 .
\]
In the representation \(\Lambda_k(u)=\sum_v\gamma_{uv}^{(k)}(\varepsilon_v)\),
all \(v\ne u\) terms have zero expectation even after multiplication by
\(g_{ktu}(\varepsilon_u)\). The remaining diagonal is bounded by
\(L_2G_\infty\alpha_{k-1}\).
This proves the score bound \(\beta_k\), including the required
first-order cancellation.

The actual population curvature at \(S_{\phi,t}^{(k)}\) is at least
\(f_{kt}-L_3\E\norm{E_k}_\infty\ge3\lambda/4\) by
\eqref{eq:target-smallness} and Jensen's inequality.
Within distance \(2\beta_k/\lambda\) it is at least \(\lambda/2\).
The population score therefore crosses zero in this interval, proving
the target bound.
\end{proof}

\begin{corollary}[All components for the paper's targets]\label{cor:pdf}
Under Theorem~\ref{thm:main} and Proposition~\ref{prop:targets}, with
probability at least \(1-e^{-x}\), for all \(k\le K\),
\[
 \norm{\widetilde S^{(k)}-S_{\rm pdf}^{(k)}}_\infty
 \le A_\Lambda\sqrt{z\alpha_k}+D(1+r)z\alpha_k+2\beta_k/\lambda,
\]
where \(\beta_1=0\).
In particular, suppose both tail terms in \eqref{eq:U}--\eqref{eq:Vmom}
are at most \(\alpha_{\min}:=\min_j\alpha_j\). Define
\[
 B_{\rm tar}=L_2G_\infty+L_2D+(L_3/2)A^2+L_2+L_3/2.
\]
Then \(\beta_k\le B_{\rm tar}(z_0+1)\alpha_{k-1}\).
Since \(z_0+1\le2z\), the bound becomes
\[
 \boxed{\quad
 \norm{\widetilde S^{(k)}-S_{\rm pdf}^{(k)}}_\infty
 \le A_\Lambda\sqrt{z\alpha_k}
     +\left\{D(1+r)+\frac{4rB_{\rm tar}}{\lambda}\right\}z\alpha_k
 \quad}
\]
simultaneously over every stage and evaluation point.
For \(\alpha_k\asymp(nh_k)^{-1}\), this is the claimed
\(\sqrt{z/(nh_k)}+z/(nh_k)\) rate with constants independent of the stage.
\end{corollary}
\begin{proof}
Use the triangle inequality and the preceding proposition.
The extra tail assumption and \(\alpha_{\min}\le\alpha_{k-1}\) give the
stated \(\beta_k\) bound directly. Apply
\(\alpha_{k-1}\le r\alpha_k\). At stage one the target correction is zero.
\end{proof}

\section{A nonempty regime with a growing number of components}
For an explicit example, let \(n\) be a power of two, \(X\) be constant,
and \(\varepsilon\sim N(0,\sigma^2)\), with \(H=2\sigma\).
At stage \(k\), average within a partition into \(2^{k-1}\) equal blocks,
each partition refining the preceding one. These weights are an example
covered by the general theorem, rather than the repository's convolution
weights. Here
\[
 \alpha_k=2^{k-1}/n,\qquad a=2,\quad C_\alpha=1,\quad V_k=W_k.
\]
The reference is constant at stage one and zero thereafter.
The averaging projections satisfy \(W_mW_l=W_lW_m=W_l\) for \(l\le m\),
so \(B_{m:l}=I-W_m\) and \(C_B=2\) suffices for all partial products.
The curvature lower bound is
\(\lambda=\sqrt{2/\pi}/\sqrt{H^2+\sigma^2}\).
With \(t=\sigma^2/H^2=1/4\), Cauchy--Schwarz gives
\[
 d\le\sqrt{\frac{\E g^2}{(\E g)^2}-1}
   =\sqrt{\frac{1+t}{\sqrt{1+2t}}-1}<\frac12
   =\frac{a-1}{C_BC_\alpha}.
\]
Thus every constant in Theorem~\ref{thm:main}, including \(\eta>0\),
can be fixed independently of \(n,K\).

Choose \(K\le c\log_2 n\) for any fixed \(0<c<1\).
Then \(z\bar\alpha=O(n^{c-1}\log n)\to0\) at fixed confidence,
and \(x_*\) is eventually at least a positive constant times \(n^{1-c}\).
The quantity \(M_4\) grows at most polynomially in \(n\); hence both
exponential tail terms are eventually smaller than
\(\alpha_{\min}=1/n\). Also \(V_k^{\rm mom}\) and \(\beta_k\) tend to zero
uniformly in \(k\), so the target-localization conditions hold.
Consequently the final all-components rate, including the second-order
target correction, holds for all sufficiently large \(n\) in this
explicit growing-stage regime. The sufficient conditions are nonvacuous.

\section{Scope and application to the submitted design}
The propagation and induction gaps are closed under
\eqref{eq:dispersion} and the explicit finest-scale restriction
\eqref{eq:eta}. The target gap is closed under the fourth-moment,
tail, and target-localization conditions above. The original theorem
without these added restrictions remains unproved.

For the submitted Gaussian example, an analytic Gaussian integration
formula gives \(d\approx0.1711181\). The finite-design operator computation
gives \(C_B\approx2.1127187\) and \(C_\alpha\approx1.0045416\), for which
the dispersion threshold is approximately \(0.1951707\). Thus the new
stability inequality passes a floating-point check in that example.
The much stronger sample-size restriction \(z\bar\alpha\le\eta\) does
not pass at \(n=1000\). The theorem therefore does not certify those
particular Monte Carlo runs. The diagnostic reports this failure explicitly.
The numerical operator values are not interval-certified constants.

For reproducibility, the Gaussian dispersion computation uses an exact
one-dimensional integration identity. Write
\(B=\sqrt{H^2+\sigma^2}\),
\(\mu(c)=\sqrt{2/\pi}B^{-1}\exp[-c^2/(2B^2)]\), and
\(t_c=H\sqrt{2\log\{(\sqrt{2/\pi}/H)/\mu(c)\}}\).
Let \(P_0(c)\) be the probability that a \(N(c,\sigma^2)\) variable lies
in \([-t_c,t_c]\), and \(P_1(c)\) the same probability for
\(N(cH^2/B^2,\sigma^2H^2/B^2)\). Completing the square gives
\[
 \E|\rho''(c+\varepsilon)-\mu(c)|
       =2\mu(c)\{P_1(c)-P_0(c)\}.
\]
Indeed, the positive part occurs exactly on \(|c+\varepsilon|\le t_c\),
the tilted Gaussian integral of \(\rho''\) there is \(\mu(c)P_1(c)\),
and the centered variable has mean zero. The diagnostic checks this
formula against adaptive quadrature.

The proof of Lemma~\ref{lem:chaos} is self-contained; sharper inequalities
are available in Gine, Latala and Zinn,
\href{https://arxiv.org/abs/math/0003228}{\emph{Exponential and moment
inequalities for U-statistics}}. Sharper constants, removal of the
dispersion restriction, and general boundary designs are separate questions.

\end{document}
